\par\noindent\textbf{(1)} 外均匀静电场下设 \(\Phi_{\rm out}=-E_0r\cos\theta+C\cos\theta/r^2\)，内场 \(\Phi_{\rm in}=-E_{\rm in}r\cos\theta\)。在 \(r=a\) 处使电势与法向位移连续，求得 \(E_{\rm in}=3E_0/(\varepsilon_r+2)\)，\(C=a^3E_0(\varepsilon_r-1)/(\varepsilon_r+2)\)。所以偶极极化率 \(\alpha=4\pi\varepsilon_0a^3(\varepsilon_r-1)/(\varepsilon_r+2)\)。
\par\noindent\textbf{(2)} 诱导偶极 \(p_0=\alpha E_0\) 随入射频率振荡。远场电场振幅 \(E_{\theta,0}=\omega^2p_0\sin\theta/(4\pi\varepsilon_0c^2r)\)。以 \(\langle S_r\rangle=\varepsilon_0cE_{\theta,0}^2/2\) 求每立体角辐射功率 \(dP/d\Omega=\omega^4|p_0|^2\sin^2\theta/(32\pi^2\varepsilon_0c^3)\)。除以入射平均强度 \(I_0=\varepsilon_0c|E_0|^2/2\)，严格得到 \(d\sigma/d\Omega=k^4|\alpha|^2\sin^2\theta/(16\pi^2\varepsilon_0^2)\)。
\par\noindent\textbf{(3)} 直接代入偶极辐射 \(dP/d\Omega=\omega^4|p_0|^2\sin^2\theta/(32\pi^2\varepsilon_0c^3)\) 及 \(I_0=\varepsilon_0c|E_0|^2/2\)，严格得到 \(\frac{d\sigma}{d\Omega}=\frac{k^4|\alpha|^2}{16\pi^2\varepsilon_0^2}\sin^2\theta\)。利用 \(\int\sin^2\theta\,d\Omega=8\pi/3\)，总散射截面 \[\sigma_{\rm sca}=k^4|\alpha|^2/(6\pi\varepsilon_0^2)=\frac{8\pi}{3}k^4a^6\left|\frac{\varepsilon_r-1}{\varepsilon_r+2}\right|^2\]。
\par\noindent\textbf{(4)} 固定波长与介电率时 \(\sigma_{\rm sca}\propto a^6\)，解释小粒径散射快速减弱；\(\varepsilon_r\to1\) 时极化率与散射消失。若 \(ka\) 不是小量或材料吸收明显，需使用复极化率、辐射反作用修正乃至Mie理论，不可将静态极化率用于任意大球。
